\input{realizations/continuation/C63/formal_immersion63}
\begin{proposition}\label{prop:nonic-no63}
An elliptic curve over a number field of degree $9$ has no point of
order $63$.
\end{proposition}
\begin{proof}
Put $Y=X_1(63)/\langle8\rangle$. The diamond $8$ acts by different
signs on the $9$- and $7$-components. This involution acts freely on
the noncuspidal locus also in characteristic $2$: an automorphism
fixing such a level structure would have square $1$, and hence
would be $1$ or $-1$, neither of which induces these two signs.
The primitive-column calculation verifies freeness at the cusps.
Thus the quotient is etale and $Y$ has good reduction at $2$.
Its genus is $49$.

Every even character modulo $63$ which is trivial at $8$ has full
winding submodule. The five character orbits have weighted dimension
sum $49$. Hence $J(Y)(\mathbb Q)$ has rank zero.
We use the direct model over $\mathbb F_2$ obtained from
\[
 E_9:y^2+(1-c)xy-by=x^3-bx^2,\qquad
 c=t^2(t-1),\quad b=t^2(t-1)(t^2-t+1),
\]
by adjoining a root $u$ of $\psi_7(E_9)$.
Here $(0,0)$ has exact order $9$, and $u=x(P_7)$ forgets the sign
of the $7$-point. The polynomial is irreducible of degree $24$.
Its function field has genus $49$, and its $j$-map has degree $864$.
The closed cusps over $\mathbb F_2$ have degrees
\[
 1^9 2^3 3^9 6.
\]
The global cusp degrees are $1^9 2^3 3^6 6\,9$.
In particular every cusp of degree at most $2$ in the special
fibre lifts to a global cusp of the same degree.
The three cubic places arising from the global degree-$9$ cusp
are retained in the calculation below.

Suppose that $D_X$ is the rational noncuspidal degree-$9$ trace
divisor associated to a point of order $63$, and let $D$ be its
pushforward to $Y$. The exact source counts at $2$ give the numbers
of noncuspidal closed places of degrees $1,\ldots,9$ as
\[
 (0,0,0,0,0,0,54,54,72).
\]
Thus a noncuspidal part in the reduction of $D_X$ has degree
$7$, $8$, or $9$, with a cuspidal complement of degree at most $2$.

We first exclude these three cases.
Direct enumeration in the displayed model gives respectively
$189$, $216$, and $324$ pairs $(t,u)$ over the fields of degrees
$7$, $8$, and $9$ which lift to actual $63$-torsion structures.
Their Frobenius orbits all have the indicated degree and give
respectively $27$, $27$, and $36$ closed places on $Y$.
They form respectively $3$, $3$, and $4$ orbits under the residual
diamonds $2,5$.
For a representative $P$ of each of these ten orbits, the exact
divisor calculation gives
\[
 A_5 P:=T_5P-P-5\langle5\rangle P
 \quad\text{nonprincipal}.
\]
The computation constructs all six $5$-isogenies, transports both
torsion points, and recovers the pair $(t,u)$ by Tate normalization.
Divisor multiplicities include the ratios of residue degrees.
For the new degree-$9$ representatives, each Hecke image is a
single place of degree $54$.

However, $A_5$ kills $J(Y)(\mathbb Q)$ by Eichler--Shimura and
injectivity of torsion reduction at the odd good prime $5$.
For a rational unit-component cusp $c_0$, we have the literal
identity $A_5c_0=0$.
Any cuspidal complement $C$ of degree at most $2$ lifts globally,
so $A_5(C-(\deg C)c_0)$ is principal.
Applying $A_5$ to the rational class $[D-9c_0]$ and then reducing
at $2$ would therefore make $A_5P$ principal. This contradicts
the ten tests.

It remains to exclude entirely cuspidal reduction.
Put $U=(\mathcal C-1)H^0(Y_{\mathbb F_2},\Omega^1)$.
The fresh Cartier computation gives $\dim U=45$.
We compute all $121$ spaces
\[
 H^0\bigl(Y_{\mathbb F_2},\Omega^1(-mP)\bigr),\qquad
 1\leq m\leq\left\lfloor\frac9{\deg P}\right\rfloor,
\]
at the $22$ closed cusps. Their annihilators give compatible
restriction matrices. Restriction from $U$ has rank $9$ at every
one of the $105040$ effective cuspidal divisors of degree $9$.
Independent binary linear algebra checks every kernel, projection,
and restriction rank.
Lemma~\ref{lem:nonic63-quotient-Cartier} now forces $D$ to be
cuspidal, a contradiction.
\end{proof}
